\documentclass{article}
\usepackage[utf8]{inputenc}
\usepackage{graphicx}
\usepackage{caption}  % Pacchetto per controllare le didascalie
\usepackage{amsmath}
\usepackage{amssymb} % Simboli matematici tipo R (ovvero numeri reali)
\title{Esercizi di Analisi}
\author{Nicolò Giuliani}
\date{}
\begin{document}
\maketitle
\section*{Esercizio 1:}
Individuare e classificare i punti critici della funzione $f : \mathbb{R}^2 \to \mathbb{R}$
\begin{equation*}
    f(x,y)=e^{x^2+y}-x^2+e^{-(y+1)}
\end{equation*}
i punti critici si trovano in $\nabla f(0,0)$:
\begin{align*}
    \frac{\partial f}{\partial x}=e^{x^2+y}*(2x+0)-2x+e^{-(y+1)}*0 \\
    = 2xe^{x^2+y}-2x \\
    = 2x(e^{x^2+y}-1)
\end{align*}
\begin{align*}
    \frac{\partial f}{\partial y}=e^{x^2+y}*(0+1)-0+e^{-(y+1)}*(-1+0) \\
    = e^{x^2+y} - e^{-(y+1)}
\end{align*}
bene, ora troviamo i punti critici:
\begin{align*}
    \frac{\partial f}{\partial x}=0 \\
    2x(e^{x^2+y}-1) = 0 \to x=0
\end{align*}
poiche' $x=0$ la seconda equazione diventa:
\begin{align*}
    e^{0^2+y} - e^{-(y+1)}=0 \\
    \to e^y - e^{-(y+1)}=0 \\
    \to e^y = e^{-y-1} \\
    \to Ln(e^y) = Ln(e^{-y-1}) \\
    \to y = -y -1 \\
    \to 2y = -1 \\
    \to y = -\frac{1}{2}
\end{align*}
Abbiamo trovato il punto critico $(0,-\frac{1}{2})$. 
\vspace{1em}
\\Ok ora troviamo gli altri punti, caso $x\neq 0$:
\begin{align*}
    2x(e^{x^2+y}-1) = 0 \\
    \to e^{x^2+y}-1=0 \\
    \to e^{x^2+y}=1 \\
    \to Ln(e^{x^2+y})=Ln(1) \\
    \to x^2+y=0 \\
    \to y = -x^2
\end{align*}
Sostituiamo $y=-x^2$ nella seconda equazione:
\begin{align*}
    e^{x^2+(-x^2)}-e^{-(-x^2+1)}=0 \\
    \to e^{0}-e^{x^2-1}=0 \\
    \to 1-e^{x^2-1}=0 \\
    \to -e^{x^2-1}=-1 \\
    \to e^{x^2-1}=1 \\
    \to Ln(e^{x^2-1})=Ln(1) \\
    \to x^2-1=0 \\
    \to x^2=1 \\
    \to x=\pm 1
\end{align*}
Ora sostituiamo a $y=-x^2$ i valori appena trovati:
\begin{align*}
    y=-(1)^2=-1 \\
    y=-(-1)^2=-1
\end{align*}
Quindi abbiamo trovato i punti critici $(1,-1)$ e $(-1,-1)$. \\
Calcoliamo ora la matrice hessiana:

\begin{align*}
    H=\begin{pmatrix}
        \frac{\partial f}{\partial^2 x} & \frac{\partial f}{\partial x \partial y} \\
        \frac{\partial f}{\partial y \partial x} & \frac{\partial f}{\partial^2 y}
    \end{pmatrix}
    = \begin{pmatrix}
        2e^{x^2+y}+4x^2e^{x^2+y}-2 & 2xe^{x^2+y} \\
        2xe^{x^2+y} & e^{x^2+y} + e^{-y-1}
    \end{pmatrix}
\end{align*}
Ok bene, ora mettiamo come valori $(0,-\frac{1}{2})$
\begin{align*}
    H_{(0,-\frac{1}{2})}=\begin{pmatrix}
        2(\frac{1}{\sqrt{e}}-1) & 0 \\
        0 & \frac{2}{\sqrt{e}}
    \end{pmatrix}\\
     \text{$det$ $H$}=2e^{0^2-\frac{1}{2}}+4*0*^2e^{0^2+\frac{1}{2}}-2(\frac{2}{\sqrt{e}}) \\
     = 2e^{-\frac{1}{2}}-\frac{4}{\sqrt{e}} \\
     = \frac{2}{\sqrt{e}}-\frac{4}{\sqrt{e}} \\
     = -\frac{2}{\sqrt{e}} < 0
\end{align*}
Ora calcoliamo la traccia
\begin{align*}
     \text{$Tr$ $H$}= 2(\frac{1}{\sqrt{e}}-1) + \frac{2}{\sqrt{e}} - 0^2=  \frac{2}{\sqrt{e}}-2+ \frac{2}{\sqrt{e}}\\     
     = \frac{4}{\sqrt{e}}-2
\end{align*}
Ricorda: $e=2,71$ quindi $\sqrt{2,71}=1,\dots > 1$ 
\begin{align*}
    = \frac{4}{1,\dots}-2 \\
    = 4,\dots -2 \\
    = 2,\dots >0
\end{align*}
Quindi se il determinanente e' negativo, la traccia non importa, quindi abbiamo un punto di sella nel punto $(0,-\frac{1}{2})$. \\ \\
Ora classifichiamo $(1,-1)$
\begin{align*}
    H_{(1,-1)}=\begin{pmatrix}
        2e^{1-1}+4e^{1-1}-2 & 2e^{1-1} \\
        2e^{1-1} & e^{1-1} + e^{+1-1}
    \end{pmatrix}
    = \begin{pmatrix}
            2e^{0}+4e^{0}-2 & 2e^{0} \\
            2e^{0} & e^{0} + e^{0}
        \end{pmatrix}
    = \begin{pmatrix}
        2+4-2 & 2 \\
        2 & 1 + 1
    \end{pmatrix} \\
    =\begin{pmatrix}
        4 & 2 \\
        2 & 2
    \end{pmatrix}
\end{align*}
\begin{align*}
    \text{$det$ $H$}= 4*2 - 2*2 = 8 - 4 = 4 >0 \\
    \text{$Tr$ $H$}=4+2=6 >0
\end{align*}
Quindi se il determinanente e' positivo e la traccia pure abbiamo un minimo locale in $(1,-1)$.


\pagebreak
Ok ci manca poco, ora facciamo $(-1,-1)$
\begin{align*}
    H_{(-1,-1)}=\begin{pmatrix}
        2e^{1-1}+4e^{1-1}-2 & -2e^{1-1} \\
        -2e^{1-1} & e^{1-1} + e^{+1-1}
    \end{pmatrix}
    = \begin{pmatrix}
            2e^{0}+4e^{0}-2 & -2e^{0} \\
            -2e^{0} & e^{0} + e^{0}
        \end{pmatrix}
    = \begin{pmatrix}
        2+4-2 & -2 \\
        -2 & 1 + 1
    \end{pmatrix} \\
    =\begin{pmatrix}
        4 & -2 \\
        -2 & 2
    \end{pmatrix}
\end{align*}
\begin{align*}
    \text{$det$ $H$}= 4*2 - (-2)(-2) = 8 - 4 = 4 >0 \\
    \text{$Tr$ $H$}=4+2=6 >0
\end{align*}
Quindi se il determinanente e' positivo e la traccia pure abbiamo un minimo locale in $(-1,-1)$.
\pagebreak
\section*{Esercizio 2:}
1) Individuare e classificare i punti critici.\\
\begin{align*}
    f(x,y)=(1+(x+1)y)^\frac{1}{2}-x^2 \\
    = \sqrt{1+(x+1)y}-x^2
\end{align*}
Calcoliamo $\nabla f(0,0)$:
\begin{align*}
    \frac{\partial f}{\partial x}=\frac{1}{2}*(1+(x+1)y)^{-\frac{1}{2}}*\frac{\partial({1+xy+y})}{\partial x}-2x =0\\
    = \frac{1}{2} * (1+(x+1)y)^{-\frac{1}{2}} * (0 + y + 0) -2x \\
    = \frac{1}{2(1+xy+y)^\frac{1}{2}}*y -2x \\
    = \frac{y}{2(1+xy+y)^\frac{1}{2}} -2x \\ 
    \to \frac{y}{2\sqrt{(1+xy+y)}} -2x =0 \\
    \to \frac{y+(-2x)(2\sqrt{(1+xy+y)})}{2\sqrt{(1+xy+y)}}=0 \\
    \to  \frac{y-4x\sqrt{(1+xy+y)}}{2\sqrt{(1+xy+y)}}=0
\end{align*}
\begin{align*}
    \frac{\partial f}{\partial y}=\frac{1}{2}*(1+(x+1)y)^{-\frac{1}{2}}*\frac{\partial(1+xy+y)}{\partial y} - 0 =0\\
    = \frac{1}{2(1+xy+y)^\frac{1}{2}}*(0+x+1) \\
    \to \frac{x+1}{2\sqrt{(1+xy+y)}} = 0  \\
\end{align*}
In $\frac{\partial f}{\partial x}$ il denominatore non e' mai nullo, quindi: 
\begin{align*}
    {y-4x\sqrt{(1+xy+y)}}=0 \\
    \to y = 4x\sqrt{(1+xy+y)}
\end{align*}
Se lo sostituiamo in $\frac{\partial f}{\partial y}$
\begin{align*}
    \frac{x+1}{2\sqrt{(1+4x^2\sqrt{(1+xy+y)}+4x\sqrt{(1+xy+y)})}} = 0 \\
    \frac{x+1}{2\sqrt{(1+4x\sqrt{(1+xy+y)}(x+1))}} = 0
\end{align*}
Anche qui il denominatore e' sempre positivo, quindi $\frac{\partial f}{\partial y}=0$ quando il numeratore e' nullo, ovvero
\begin{align*}
    x+1=0 \\
    \to x=-1
\end{align*}
Sostituiamo $x=-1$ in $\frac{\partial f}{\partial x}$
\begin{align*}
    \frac{y}{2(1-y+y)^\frac{1}{2}} +2 =0\\ 
    = \frac{y}{2(\sqrt{1})}+2 \\
    = \frac{y}{2}+2 =0\\
    \to \frac{y}{2}=-2 \\
    \to y = -4
\end{align*}
Quindi $(-1,-4)$ e' un punto critico. \\
Anche in $\frac{\partial f}{\partial y}$ notiamo che il denominatore non e' mai negativo, quindi per avere l'equazione uguale a 0 dobbiamo avere il numeratore nullo, quindi:
\begin{align*}
    x+1=0 \\
    x=-1
\end{align*}
Anche qui abbiamo ritrovato $x=-1$ quindi $(-1,-4)$ e' l'unico punto critico. \\
Bene ora calcoliamo le derivate seconde per la matrice hessiana:
\begin{align*}
    \frac{\partial f}{\partial^2 x}=\frac{\partial(\frac{y}{2\sqrt{(1+xy+y)}})}{\partial x}-2 =0\\
    = -2 + \frac{\frac{\partial y}{\partial x}* 2\sqrt{(1+xy+y)} - \frac{\partial (2\sqrt{(1+xy+y)})}{\partial x}*y}{(2\sqrt{(1+xy+y)}^2)} \\
    = -2 - \frac{y(2*(-\frac{1}{2})*(1+xy+y)^{-\frac{1}{2}}*\frac{\partial (1+xy+y)}{\partial x})}{4(1+xy+y)} \\
    = -2 - \frac{y(-1*(1+xy+y)^{\frac{1}{2}} * (y))}{4(1+xy+y)} \\
    = -2 - \frac{y^2 * (1+xy+y)^{-\frac{1}{2}}}{4(1+xy+y)} \\
    = -2 - \frac{y^2}{4(1+xy+y)*(1+xy+y)^{\frac{1}{2}}} \\
    = -2 - \frac{y^2}{4(1+xy+y)\sqrt{(1+xy+y)}} \\
    = -2 - \frac{y^2}{4(1+xy+y)^{\frac{3}{2}}} \\
    = -2 - \frac{y^2}{4\sqrt{(1+xy+y)^3}} =0
\end{align*}
\begin{align*}
    \frac{\partial f}{\partial x \partial y}=\frac{\frac{\partial y}{\partial y}*(2(1+xy+y)^{\frac{1}{2}})-\frac{\partial (2(1+xy+y)^{\frac{1}{2}})}{\partial y}*y}{(2(1+xy+y)^{\frac{1}{2}})^2} - 0 =0\\
    = \frac{(2(1+xy+y)^{\frac{1}{2}}) - y(2*\frac{1}{2}*(1+xy+y)^{-\frac{1}{2}}*\frac{\partial (1+xy+y)}{\partial y})}{4(1+xy+y)} \\
    = \frac{(2(1+xy+y)^{\frac{1}{2}}) - y((1+xy+y)^{-\frac{1}{2}}*(x+1))}{4(1+xy+y)} \\
    = \frac{(2\sqrt{(1+xy+y)}) - (xy+y)(\sqrt{(1+xy+y)}^{-1})}{4(1+xy+y)} \\
    = \frac{2\sqrt{(1+xy+y)}}{4(1+xy+y)} - \frac{xy+y}{4(1+xy+y)(1+xy+y)^{\frac{1}{2}}} \\
    = \frac{\sqrt{(1+xy+y)}}{2(1+xy+y)} - \frac{xy+y}{4(1+xy+y)^{\frac{3}{2}}} \\
    = \frac{\sqrt{(1+xy+y)}(2(1+xy+y)^{\frac{1}{2}})}{4(1+xy+y)^{\frac{3}{2}}} - \frac{xy+y}{4(1+xy+y)^{\frac{3}{2}}} \\
    = \frac{2(1+xy+y)}{4(1+xy+y)^{\frac{3}{2}}} -\frac{xy+y}{4(1+xy+y)^{\frac{3}{2}}} \\
    = \frac{2+2xy+2y-xy-y}{4(1+xy+y)^{\frac{3}{2}}} \\
    = \frac{2+xy+y}{4(1+xy+y)^{\frac{3}{2}}}=0
\end{align*}
\begin{align*}
    \frac{\partial f}{\partial y\partial x}=\frac{\frac{\partial (x+1)}{\partial x}*2(1+xy+y)^{\frac{1}{2}}-(x+1)(\frac{\partial 2(1+xy+y)^{\frac{1}{2}}}{\partial x})}{(2(1+xy+y)^{\frac{1}{2}})^2} =0\\
    = \frac{2(1+xy+y)^{\frac{1}{2}} - (x+1)*2*\frac{1}{2}*(1+xy+y)^{-\frac{1}{2}}*\frac{\partial (xy+y)}{\partial x}}{4(1+xy+y)} \\
    = \frac{2(1+xy+y)^{\frac{1}{2}} - (x+1)*(1+xy+y)^{-\frac{1}{2}}* y}{4(1+xy+y)} \\
    = \frac{2(1+xy+y)^{\frac{1}{2}} - (yx+y)*(1+xy+y)^{-\frac{1}{2}}}{4(1+xy+y)} \\
    = \frac{2(1+xy+y)^{\frac{1}{2}}}{4(1+xy+y)} - \frac{yx+y}{4(1+xy+y)(1+xy+y)^{\frac{1}{2}}} \\
    = \frac{2(1+xy+y)^{\frac{1}{2}}}{4(1+xy+y)} - \frac{yx+y}{4(1+xy+y)^{\frac{3}{2}}} \\
    = \frac{(1+xy+y)^{\frac{1}{2}}}{2(1+xy+y)} - \frac{yx+y}{4(1+xy+y)^{\frac{3}{2}}} \\
    = \frac{(1+xy+y)^{\frac{1}{2}}(2(1+xy+y)^{\frac{1}{2}})}{4(1+xy+y)^{\frac{3}{2}}} - \frac{yx+y}{4(1+xy+y)^{\frac{3}{2}}} \\
    = \frac{2(1+xy+y)}{4(1+xy+y)^{\frac{3}{2}}} - \frac{yx+y}{4(1+xy+y)^{\frac{3}{2}}} \\
    = \frac{2+2xy+2y -xy -y}{4(1+xy+y)^{\frac{3}{2}}} \\
    = \frac{2+xy+y}{4(1+xy+y)^{\frac{3}{2}}}=0 
\end{align*}
\begin{align*}
    \frac{\partial f}{\partial^2 y}=-\frac{(x+1)* \frac{\partial (2(1+xy+y)^{\frac{1}{2}})}{\partial y}}{(2(1+xy+y)^{\frac{1}{2}})^2} =0\\
    =-\frac{(x+1)* 2*\frac{1}{2}*(1+xy+y)^{-\frac{1}{2}}*\frac{\partial (xy+y)}{\partial x}}{4(1+xy+y)} \\
    =-\frac{(x+1)*(1+xy+y)^{-\frac{1}{2}} * y}{4(1+xy+y)} \\
    =-\frac{(yx+y)*(1+xy+y)^{-\frac{1}{2}}}{4(1+xy+y)} \\
    =-\frac{yx+y}{4(1+xy+y)*(1+xy+y)^{\frac{1}{2}}} \\
    =-\frac{yx+y}{4(1+xy+y)^{\frac{3}{2}}}=0
\end{align*}

Scriviamo la matrice:
\begin{align*}
    H = \begin{pmatrix}
        -2 - \frac{y^2}{4\sqrt{(1+xy+y)^3}} & \frac{2+3xy+3y}{4(1+xy+y)^{\frac{3}{2}}} \\
        \frac{2+3xy+3y}{4(1+xy+y)^{\frac{3}{2}}} & -\frac{yx+y}{4(1+xy+y)^{\frac{3}{2}}}
    \end{pmatrix}
\end{align*}
Ora calcoliamo per il punto $(-1,-4)$
\begin{align*}
    H_{(-1,-4)} = \begin{pmatrix}
        -2 - \frac{(-4)^2}{4\sqrt{(1+(-1)(-4)-4)^3}} & \frac{2+3(-1)(-4)+3(-4)}{4(1+(-1)(-4)-4)^{\frac{3}{2}}} \\
        \frac{2+3(-1)(-4)+3(-4)}{4(1+(-1)(-4)-4)^{\frac{3}{2}}} & -\frac{(-1)(-4)-4}{4(1+(-1)(-4)-4)^{\frac{3}{2}}}
    \end{pmatrix} \\
    =\begin{pmatrix}
        -2 - \frac{16}{4\sqrt{(1)^3}} & \frac{2}{4(1)^{\frac{3}{2}}} \\
        \frac{2}{4(1)^{\frac{3}{2}}} & -\frac{0}{4(1)^{\frac{3}{2}}}
    \end{pmatrix} 
    = \begin{pmatrix}
        -2 - \frac{16}{4} & \frac{2}{4} \\
        \frac{2}{4} & 0
    \end{pmatrix} 
    = \begin{pmatrix}
        -2 - 4 & \frac{1}{2} \\
        \frac{1}{2} & 0
    \end{pmatrix} \\
    = \begin{pmatrix}
        -6 & \frac{1}{2} \\
        \frac{1}{2} & 0
    \end{pmatrix}
\end{align*}

Calcoliamo ora il determinanente
\begin{align*}
    \text{$det$ $H$}= 2 * 0 - (\frac{1}{2})^2\\
    = -\frac{1}{4} <0
\end{align*}
Dato che il determinante nel punto $(-1,-4)$ e' negativo abbiamo un punto di sella.

\pagebreak
\section*{Esercizio 3:}
Calcolare nell'insieme $A=\{(x,y) \in \mathbb{R}^2 : |x|\leq\pi , y \geq \frac{1}{2} \text{ e } y\leq sin(x)\}$ l'integrale:
\begin{equation*}
    \int_A y \,dxdy
\end{equation*}
Sappiamo che A e' limitato da
\begin{align*}
   \text{Orizzontalmente: } -\pi \leq x \leq \pi \\
    \text{Verticalmente: } \frac{1}{2}\leq y \leq sin(x) 
\end{align*}
Riscriviamo l'integrale come
\begin{align*}
    \int_{-\pi}^\pi \int_{\frac{1}{2}}^{sin(x)} y \,dydx
\end{align*}
Calcoliamo il primo integrale
\begin{align*}
    \int_{\frac{1}{2}}^{sin(x)} y \,dy=[\frac{y^2}{2}]_{\frac{1}{2}}^{sin(x)} \\
    = \frac{sin^2(x)}{2}-\frac{\frac{1}{2}^2}{2} \\
    = \frac{sin^2(x)}{2}-\frac{1}{8}
\end{align*}
Bene ora inseriamolo nel secondo integrale
\begin{align*}
    \int_{-\pi}^{\pi} \frac{sin^2(x)}{2}-\frac{1}{8} \,dx
\end{align*}
ovvero
\begin{align*}
    \int_{-\pi}^{\pi} \frac{sin^2(x)}{2} \,dx + \int_{-\pi}^{\pi} -\frac{1}{8} \,dx \\
    = \int_{-\pi}^{\pi} \frac{sin^2(x)}{2} \,dx -\frac{1}{8} \int_{-\pi}^{\pi}  \,dx
\end{align*}
Calcoliamo il primo integrale:
\begin{align*}
    \int_{-\pi}^{\pi} \frac{sin^2(x)}{2} \,dx= \frac{1}{2}\int_{-\pi}^{\pi} sin^2(x) \,dx \\
\end{align*}
Ricorda la formula di riduzione:
\begin{align*}
    sin^2(x)=\frac{1-cos(2x)}{2}
\end{align*}
Quindi
\begin{align*}
    \frac{1}{2}\int_{-\pi}^{\pi} \frac{1-cos(2x)}{2} \,dx \\
    = \frac{1}{4}\int_{-\pi}^{\pi} 1-cos(2x) \,dx \\
    = \frac{1}{4}(\int_{-\pi}^{\pi} 1 \,dx - \int_{-\pi}^{\pi} cos(2x) \,dx) \\
    = \frac{1}{4}([x]_{-\pi}^{\pi} - \int_{-\pi}^{\pi} cos(2x) \,dx)
\end{align*}
Ricorda:
\begin{align*}
    \int cos(2x) \, dx = \frac{1}{2}\int 2cos(2x) \,dx \\
    \text{Sappiamo che: } \int cos(\alpha x) \, dx = \frac{sin(\alpha x)}{\alpha} \\
    \to \int cos(2x) \, dx = \frac{sin(2x)}{2}+C
\end{align*}
Quindi sapendo questo
\begin{align*}
    \frac{1}{4}([x]_{-\pi}^{\pi} - \int_{-\pi}^{\pi} cos(2x) \,dx) \\ 
    = \frac{1}{4}([x]_{-\pi}^{\pi} - [\frac{sin(2x)}{2}_{-\pi}^{\pi}]) \\
    = \frac{1}{4}( [x-\frac{sin(2x)}{2}_{-\pi}^{\pi}]) \\
    = \frac{1}{4}(\pi-(-\pi) + (-\frac{sin(2\pi)}{2})-(-\frac{sin(-2\pi)}{2})) \\
    = \frac{1}{4}(2\pi + 0) \\
    = \frac{2\pi}{4} \\
    = \frac{\pi}{2}
\end{align*}
Calcoliamo il secondo integrale:
\begin{align*}
    -\frac{1}{8}\int_{-\pi}^{\pi}dx \\
    = -\frac{1}{8}[x]_{-\pi}^{\pi} \\
    = -\frac{1}{8}(\pi - (-\pi)) \\
    = -\frac{1}{8}(2\pi) \\
    = -\frac{2\pi}{8} \\
    = -\frac{\pi}{4}
\end{align*}
Quindi inseriamo i valori trovati
\begin{align*}
    \int_{-\pi}^{\pi} \frac{sin^2(x)}{2} \,dx -\frac{1}{8} \int_{-\pi}^{\pi}  \,dx \\
    = \frac{\pi}{2} - \frac{\pi}{4} \\
    = \frac{2\pi - \pi}{4}\\
    = \frac{\pi}{4}
\end{align*}

\pagebreak
\section*{Esercizio 4:}
\textit{Non sono sicuro che lo svolgimento sia giusto} \\
Dato l'insieme $A={(x,y)\in \mathbb{R}^2 : x \in (0,1) \text{ e } 0\leq y \leq \sqrt{x}}$ calcolare l'integrale
\begin{equation*}
    \int_A \frac{y}{x+y^2+1} \,dxdy
\end{equation*}
Sappiamo che A e’ limitato da
\begin{align*}
    0 < x < 1 \\
    0 \leq y \leq \sqrt{x} \to x\geq y^2
\end{align*}
Riscriviamo l'integrale
\begin{align*}
    \int_0^1 \int_0^{\sqrt{x}} \frac{y}{x+y^2+1} \,dydx
\end{align*}
Cambio dell'ordine di integrazione
\begin{align*}
    0 \leq y \leq \sqrt{x} \text{ se } x=1 \to 0 \leq y \leq 1 \\
    0 \leq x \leq 1 \text{ e } x\geq y^2 \leq 1 \to y^2 \leq x \leq 1 \\
    \int_0^{1}\int_{y^2}^1  \frac{y}{x+y^2+1} \,dxdy
\end{align*}
Integrazione rispetto a $x$
\begin{align*}
    \int_{y^2}^1  \frac{y}{x+y^2+1} \,dx
\end{align*}
Integriamo per sostituzione
\begin{align*}
    u=x+y^2+1 \\
    \frac{du}{dx}=\frac{d(x+y^2+1)}{dx}=1 \to du=dx
\end{align*}
Quindi se sostituiamo dentro l'integrale
\begin{align*}
    \int_{y^2}^1  \frac{y}{u} \,du \\
    = y \int_{y^2}^1 u \,du \\
    = y[Ln(u)]_{y^2}^1 \\
    = y[Ln(x+y^2+1)]_{y^2}^1 \\
    =y(Ln(y^2+2)-Ln(2y^2+1))
\end{align*}
Integrazione rispetto a $y$
\begin{align*}
    \int_0^1 y(Ln(y^2+2)-Ln(2y^2+1)) \,dy \\
    =  \int_0^1 y(Ln(y^2+2)) \,dy - \int_0^1 y(Ln(2y^2+1)) \,dy
\end{align*}
Integriamo per sotstituzione il primo integrale
\begin{align*}
    u=y^2+2 \\
    du=2ydy \iff dy=\frac{du}{2y}
\end{align*}
Se $y=2 \to u=0^2+2=2$ \\
Se $y=1 \to u=1^2+2=3$
\begin{align*}
    \int_2^3 y(Ln(u))\frac{du}{2y} \\
    = \int_2^3 (Ln(u))\frac{du}{2} \\
    = \frac{1}{2}\int_2^3 Ln(u) \,du
\end{align*}
Ricorda:
\begin{align*}
    \int Ln(u) \,du = uLn(u)-u +C
\end{align*}
Quindi
\begin{align*}
    \frac{1}{2}[uLn(u)-u]_2^3 \\
    = \frac{1}{2}(3Ln(3)-3 - (2Ln(2)-2)) \\
    = \frac{1}{2}(3Ln(3)-2Ln(2)-3+2) \\
    = \frac{1}{2}(3Ln(3)-2Ln(2)-1) 
\end{align*}
Integriamo per sotstituzione il secondo integrale
\begin{align*}
    - \int_0^1 y(Ln(2y^2+1)) \,dy \\
    u=2y^2+1 \\
    du=4y \,dy \iff dy=\frac{du}{4y} \\
    \text{Per y=0} \to u=2(0)^2+1=1 \\
    \text{Per y=1} \to u=2(1)^2+1=3 \\
    = - \int_1^3 y(Ln(u)) \,\frac{du}{4y} \\
    = - \int_1^3 Ln(u) \,\frac{du}{4} \\
    = -\frac{1}{4} \int_1^3 Ln(u) \,du \\
    = - \frac{1}{4}[uLn(u)-u]_1^3 \\
    = -\frac{1}{4}(3Ln(3)-3 - (Ln(1)-1)) \\
    = -\frac{1}{4}(3Ln(3)-3 - Ln(1)+1) \\
    \text{Ricorda la proprieta' dei logaritmi: } \alpha Ln(y)=Ln(y^\alpha) \\
    = -\frac{1}{4}(Ln(9)-3 - 0 +1) \\
    = -\frac{1}{4}(Ln(9)-2)
\end{align*}
Quindi
\begin{align*}
    \int_0^1 y(Ln(y^2+2)) \,dy - \int_0^1 y(Ln(2y^2+1)) \,dy \\
    = \frac{1}{2}(3Ln(3)-2Ln(2)-1) -\frac{1}{4}(Ln(9)-2) \\
    = \frac{1}{2}(Ln(3^3)-Ln(2^2)-1) -\frac{1}{4}(Ln(9)-2)\\
    = \frac{1}{2}(Ln(27)-Ln(4)-1) -\frac{1}{4}(Ln(9)-2)\\
    \text{Ricorda la proprieta' dei logaritmi: } Ln(\alpha)-Ln(\beta)=Ln(\frac{\alpha}{\beta}) \\
    = \frac{1}{2}(Ln(\frac{27}{4}-1)) -\frac{1}{4}(Ln(9)-2)\\
    \text{Facolativo: } 1=Ln(e) -\frac{1}{4}(Ln(9)-2)\\
    = \frac{1}{2}(Ln(\frac{27}{4}-Ln(e))) -\frac{1}{4}(Ln(9)-2)\\
    = \frac{1}{2}(Ln(\frac{27}{4e})) -\frac{1}{4}(Ln(9)-2)\\
    = Ln((\frac{27}{4e})^\frac{1}{3}) \-\frac{1}{4}(Ln(9)-2)\
    = Ln(\sqrt[2]{\frac{27}{4e}}) -\frac{1}{4}(Ln(9)-2)\\
    = Ln(\frac{3\sqrt{3}}{2\sqrt{e}})-\frac{1}{4}(Ln(9)-2) \\
    = Ln(3\sqrt{3})-Ln(2\sqrt{e})-\frac{1}{4}(Ln(9))+\frac{2}{4} \\
    = Ln(3)+Ln(\sqrt{3})-Ln(2)-Ln(\sqrt{e})-\frac{1}{4}(Ln(9))+\frac{1}{2} \\
    = Ln(e)+\frac{1}{2}Ln(3)-Ln(2)-\frac{1}{2}Ln(e)-\frac{1}{3}Ln(9)+\frac{1}{2} \\
    = 1 + \frac{1}{3}Ln(3)-Ln(2)-\frac{1}{2}-\frac{1}{3}Ln(9)+\frac{1}{2} \\
    = \frac{1}{3}Ln(3)-Ln(2)-\frac{1}{3}Ln(9)+1\\
    = \frac{1}{3}(Ln(3)-Ln(9))-Ln(2)+1 \\
    = \frac{1}{3}(Ln(\frac{3}{9}))-Ln(2)+1 \\
\end{align*}
\begin{align*}
    = \frac{1}{3}(Ln(\frac{1}{3}))-Ln(2)+1 \\
    = \frac{1}{3}(Ln(1)-Ln(3))-Ln(2)+1 \\
    = \frac{1}{3}(-Ln(3))-Ln(2)+1\\
    = -\frac{1}{3}Ln(3)-Ln(2)+1 \\
    = -Ln(3^\frac{1}{3})-Ln(2)+1 \\
    = -Ln(\sqrt[3]{3})-Ln(2)+1 \\
    = -Ln(\frac{\sqrt[3]{3}}{2})+1 \\
    \text{Ricorda: }-Ln(\frac{\alpha}{1})=Ln(\frac{1}{\alpha}) \\
    = Ln(\frac{2}{\sqrt[3]{3}})+1 \\
    = Ln(\frac{2}{\sqrt[3]{3}})+Ln(e) \\
    = Ln(\frac{2e}{\sqrt[3]{3}})
\end{align*}

\pagebreak
\section*{Esercizio 5:}
Dato l'insieme $A={(x,y)\in\mathbb{R}^2 : 1-x\leq y\leq 1+\frac{x}{3}\leq \frac{4}{3}}$ calcolare l'integrale
\begin{align*}
    \int_A \frac{x}{\sqrt{y+x}} \,dxdy
\end{align*}
Determiniamo i limiti di integrazione
\begin{align*}
    \text{Verticalmente limite inferiore : } 1-x\leq y \\
    \text{Verticalmente limite superiore : } y\leq 1+\frac{x}{3} \text{ e } y \leq \frac{4}{3} \to y \leq min(1+\frac{x}{3},\frac{4}{3})
\end{align*}
ovvero 
\begin{align*}
    \int \int_{1-x}^{1+\frac{x}{3}} \frac{x}{\sqrt{y+x}} \,dydx
\end{align*}
\textit{(i limiti di x non gli abbiamo ancora definiti)} \\
Determiniamo il valore di x per il limite superiore
\begin{align*}
    1+\frac{x}{3}=\frac{4}{3} \\
    \frac{x}{3}=\frac{4}{3}-1 \\
    \frac{x}{3}=\frac{1}{3} \\
    x=1
\end{align*}
Quindi abbiamo 2 casi:
\begin{align*}
    \text{1) } x \geq 1 \to y\leq\frac{4}{3} \\
    \text{2) } x \leq 1 \to y\leq 1+\frac{x}{3}
\end{align*}
Caso $x\geq 1$: \\
Dobbiamo pero' anche trovare il limite inferiore di x
\begin{align*}
    1-x\leq 1 \\
    -x\leq 1-1 \\
    x\geq 0
\end{align*}
\textit{Questo vale anche per $x\leq 1$.} \\
Quindi riscriviamo l'integrale con i nuovi limiti con $x\geq1 \to x\in (1,+\infty)$
\begin{align*}
    \int_1^{+\infty} \int_{1-x}^{\frac{4}{3}} \frac{x}{\sqrt{y+x}} \,dydx
\end{align*}
Caso $x\leq 1$: \\
Scriviamo l'integrale con $x\leq 1 \to x \in (0,1)$
\begin{align*}
    \int_0^{1} \int_{1-x}^{1+\frac{x}{3}} \frac{x}{\sqrt{y+x}} \,dydx
\end{align*}
Calcoliamo l'integrale indefinito su $y$:
\begin{align*}
    \int \frac{x}{\sqrt{y+x}} \,dy
\end{align*}
Integriamo per sostituzione
\begin{align*}
    u=(y+x)^{\frac{1}{2}} \\
    du=(\frac{1}{2}*(y+x)^{-\frac{1}{2}}*(1+0))dy \\
    = (\frac{1}{2(y+x)^{\frac{1}{2}}})dy \iff dy= du*(2(y+x)^{\frac{1}{2}}) \\
    = (\frac{1}{2(y+x)^{\frac{1}{2}}})dy \iff dy= du*(2u)
\end{align*}
Quindi l'integrale diventa
\begin{align*}
    \int \frac{x}{u} \,(2u) \,du \\
    = \int \frac{x(2u)}{u} \, \,du \\
    = \int 2x \, du \\
    = 2x \int \,du \\
    = 2x u +C \\
    = 2x (y+x)^{\frac{1}{2}} +C\\
    = 2x \sqrt{y+x}+C
\end{align*}
Calcoliamo ora per il caso $x\in(1,+\infty)$:
\begin{align*}
    \int_{1-x}^{\frac{4}{3}} \frac{x}{\sqrt{y+x}} \,dy = [2x \sqrt{y+x}]_{1-x}^{\frac{4}{3}} \\
    = (2x \sqrt{\frac{4}{3}+x} - 2x \sqrt{1-x+x}) \\
    = (2x \sqrt{\frac{4}{3}+x} - 2x \sqrt{1}) \\
    = (2x \sqrt{\frac{4}{3}+x} - 2x) 
\end{align*}
Integriamo per $x$:
\begin{align*}
    \int_1^{+\infty} 2x \sqrt{\frac{4}{3}+x} - 2x \, dx \\
    = \int_1^{+\infty} 2x \sqrt{\frac{4}{3}+x}\, dx - \int_1^{+\infty} 2x \, dx
\end{align*}
Calcoliamo il primo integrale:
\begin{align*}
    \int_1^{+\infty} 2x \sqrt{\frac{4}{3}+x}\, dx \\
    u=\frac{4}{3}+x \iff x=u-\frac{4}{3}\\
    du=dx \\
    = \int_1^{+\infty} 2x \sqrt{u} \,du  \\
    = \int_1^{+\infty} 2(u-\frac{4}{3}) \sqrt{u} \,du \\
    = \int_1^{+\infty} (2u-\frac{8}{3}) \sqrt{u} \,du
    \\\text{(i limiti sono ancora riferiti a $x$)}\\
    \text{Calcoliamo i limiti per $u$: } \\
    \text{se $x=1\to$ } u=\frac{4}{3}+1=\frac{7}{3} \\
    \text{se $x=+\infty \to$ } u=+\infty \\
    \text{Quindi con i nuovi limiti: } \\
    = \int_{\frac{7}{3}}^{+\infty} (2u-\frac{8}{3})  \sqrt{u} \,du \\
    = \int_{\frac{7}{3}}^{+\infty} 2u\sqrt{u}\,du - \frac{8}{3}\int_{\frac{7}{3}}^{+\infty}\sqrt{u}\,du
\end{align*}
Dividiamo i calcoli:
\begin{align*}
    \int_{\frac{7}{3}}^{+\infty} 2u\sqrt{u}\,du \\
    =2\int_{\frac{7}{3}}^{+\infty} u\sqrt{u}\,du \\
    =2\int_{\frac{7}{3}}^{+\infty} u^{\frac{3}{2}}\,du \\
    =2[\frac{u^{\frac{5}{2}}}{\frac{5}{2}}] \\
    =2[\frac{2u^{\frac{5}{2}}}{5}]_{\frac{7}{3}}^{+\infty} \\
    =2[\frac{2\sqrt{u^5}}{5}]_{\frac{7}{3}}^{+\infty} \\
    =2[\frac{2u^2\sqrt{u}}{5}]_{\frac{7}{3}}^{+\infty} \\
    = 2(\frac{2(+\infty)^2\sqrt{\infty}}{5} - \frac{2(\frac{7}{3})^2\sqrt{\frac{7}{3}}}{5}) \\
    = +\infty
\end{align*}
\begin{align*}
    - \frac{8}{3}\int_{\frac{7}{3}}^{+\infty}\sqrt{u}\,du\\
    \vdots \\
    =-\frac{8}{3}[\frac{2}{3}\sqrt{x^3}]_{\frac{7}{3}}^{+\infty} \\
    =-\frac{8}{3}(\frac{2}{3}\sqrt{(+\infty)^3} - \frac{2}{3}\sqrt{(\frac{7}{3})^3}) \\
    =-\frac{8}{3}(+\infty) \\
    =-\infty
\end{align*}
Ok rimettiamo tutto insieme:
\begin{align*}
    \int_{\frac{7}{3}}^{+\infty} 2u\sqrt{u}\,du - \frac{8}{3}\int_{\frac{7}{3}}^{+\infty}\sqrt{u}\,du\\
    = +\infty -\infty \\
    \text{Risultato indeterminato}
\end{align*}
Quindi:
\begin{align*}
    = \int_1^{+\infty} 2x \sqrt{\frac{4}{3}+x}\, dx - \int_1^{+\infty} 2x \, dx \\
    \text{Non si puo' calcolare perche' indeterminato}
\end{align*}
Calcoliamo ora per il caso $x\in(0,1)$
\begin{align*}
    \int_0^{1} \int_{1-x}^{1+\frac{x}{3}} \frac{x}{\sqrt{y+x}} \,dydx
\end{align*}
Integriamo prima per $y$:
\begin{align*}
    \int_{1-x}^{1+\frac{x}{3}}\frac{x}{\sqrt{y+x}} \,dy = [2x \sqrt{y+x}]_{1-x}^{1+\frac{x}{3}} \\
    = 2x \sqrt{(1+\frac{x}{3})+x} - 2x \sqrt{(1-x)+x} \\
    = 2x \sqrt{1+\frac{4x}{3}}-2x\sqrt{1} \\
    = 2x(\sqrt{1+\frac{4x}{3}}-1)
\end{align*}
Integriamo ora per $x$:
\begin{align*}
    \int_0^{1} 2x(\sqrt{1+\frac{4x}{3}}-1)\,dx \\
    =2\int_0^{1} x(\sqrt{1+\frac{4x}{3}}-1)\,dx \\
    =2\int_0^{1} x\sqrt{1+\frac{4x}{3}}-x\,dx \\
    =2(\int_0^{1} x\sqrt{1+\frac{4x}{3}}\,dx-\int_0^{1} x\,dx) \\
\end{align*}
Dividiamo i calcoli:
\begin{align*}
    \int_0^{1} x\sqrt{1+\frac{4x}{3}}\,dx \\
    u=1+\frac{4x}{3} \iff \frac{4x}{3}=u-1 \to x=(u-1)\frac{3}{4}=\frac{3u}{4}-\frac{3}{4}\\
    du=\frac{3}{4}dx \iff dx=\frac{4}{3}du \\
    =\int_0^{1} x(u)^{\frac{1}{2}} \frac{4}{3}du \\
    =\int_0^{1} ((u-1)\frac{3}{4})(u)^{\frac{1}{2}} \frac{4}{3}du \\
    =\int_0^{1} (u-1)(u)^{\frac{1}{2}} du \\
    =\int_0^{1} (u)^{\frac{3}{2}} \,du -\int_0^{1} u^{\frac{1}{2}}\,du\\
    \textit{I limiti degli integrali sono ancora su $x$} \\
    \text{Se $x=0 \to$ } u=1+\frac{4*0}{3}=1 \\
    \text{Se $x=1 \to$ } u=1+\frac{4*1}{3}=\frac{7}{3}\\
    =\int_{1}^{\frac{7}{3}} (u)^{\frac{3}{2}} \,du -\int_{1}^{\frac{7}{3}}u^{\frac{1}{2}}\,du\\
    =[\frac{u^{\frac{5}{2}}}{\frac{5}{2}}]_{1}^{\frac{7}{3}} - [\frac{2}{3}\sqrt{u^3}]_{1}^{\frac{7}{3}} \\
    =[\frac{2\sqrt{u^5}}{5}]_{1}^{\frac{7}{3}}- [\frac{2}{3}\sqrt{u^3}]_{1}^{\frac{7}{3}} \\
    =[\frac{2u^2\sqrt{u}}{5}]_{1}^{\frac{7}{3}}- [\frac{2}{3}u\sqrt{u}]_{1}^{\frac{7}{3}} \\
    = (2(\frac{7}{3})^2 \sqrt{\frac{7}{3}} - 2(1)^2\sqrt{1} ) - (\frac{2}{3}(\frac{7}{3})\sqrt{\frac{7}{3}} - \frac{2}{3}(1)\sqrt{1}) \\
    = (2(\frac{49}{9})\sqrt{\frac{7}{3}} -2)-(\frac{14}{9}\sqrt{\frac{7}{3}} -\frac{2}{3}) \\
    = \sqrt{\frac{7}{3}}(2\frac{49}{9} - \frac{14}{9})-2+\frac{2}{3} \\
    = \sqrt{\frac{7}{3}}(\frac{98}{9} - \frac{14}{9})-\frac{4}{3} \\
    = \sqrt{\frac{7}{3}}(\frac{84}{9} )-\frac{4}{3} \\
    = \frac{28}{3} \sqrt{\frac{7}{3}}-\frac{4}{3} \\
\end{align*}
\section*{Formula vettoriale del piano tangente}
Il piano tangente al grafico di $f(x,y)f(x,y)$ nel punto $(x_0,y_0)$ si scrive:
\begin{align*}
    =f(x_0,y_0)+\langle \nabla f(x_0,y_0),(x-x_0,y-y_0) \rangle
\end{align*}
Dove: \\
$\nabla f(x_0,y_0)=(f_x(x_0,y_0),f_y(x_0,y_0))$ e' il gradiente. \\
$\langle \vec{.} \, , \, \vec{.} \rangle$ e' il prodotto scalare tra due vettori in $\mathbb{R}^2$. \\
$(x-x_0,y-y_0)$ è il vettore di spostamento nel piano.
\end{document}